\section{The physical boundary experiment}\label{sec:boundary}
The algebraic extension through $v=0$ in Section~\ref{sec:native} is useful for differentiation, but it does not remove the physical restriction $v\ge0$. We now identify the actual local experiment. The argument is the finite-cell likelihood expansion underlying constrained Gaussian limits \cite{Geyer,vdV}; all normalizations and the synchronized profile are specified here.

Let $u\in\R^{k+1}$ denote the nuisance log-odds and root-sum coordinates. Observe independent categorical samples with $n_j/N\to\lambda_j>0$, and let $\Delta_N=(\Delta_{N,u},\Delta_{N,v})$ be the score at $(u,v)=(0,0)$ divided by $\sqrt N$. Partition
\[
 I=V^{\mathsf T}\diag(\lambda_j\kappa_j)V
   =\begin{pmatrix}I_{uu}&I_{uv}\\I_{vu}&I_{vv}\end{pmatrix},
 \qquad Q=I_{vv}-I_{vu}I_{uu}^{-1}I_{uv}=R^{-1}.
\]
\begin{theorem}[Cone likelihood and synchronized profile]\label{thm:boundary}
Fix a compact set $K\subset\R^{k+1}\times\R^k_{\ge0}$, with the Euclidean product norm. For $(a,h)\in K$, the physical local alternatives $u=a/\sqrt N$, $v=h/\sqrt N$ satisfy, uniformly on compact index sets,
\begin{equation}\label{eq:LAN}
 \log\frac{dP_{N,a,h}}{dP_{N,0,0}}
  =(a,h)^{\mathsf T}\Delta_N-\tfrac12(a,h)^{\mathsf T}I(a,h)+o_P(1).
\end{equation}
Here uniformity means that the supremum over $K$ of the absolute remainder converges to zero in $P_{N,0,0}$-probability. Under the corresponding alternatives,
\begin{equation}\label{eq:efficient-Y}
 Y_N=R\bigl(\Delta_{N,v}-I_{vu}I_{uu}^{-1}\Delta_{N,u}\bigr)
          \ \Longrightarrow\ N(h,R).
\end{equation}
After quotienting out unrestricted nuisance translation, this is a Gaussian shift with parameter restricted to the physical cone. For synchronized alternatives $h=q_A(\zeta)$, the profiled limiting log likelihood, relative to its value at $h=0$, is
\begin{equation}\label{eq:profile-limit}
 \frac12Y^{\mathsf T}R^{-1}Y-\frac12C_R(Y),\qquad
 C_R(Y)=\inf_\zeta(Y-q_A(\zeta))^{\mathsf T}R^{-1}(Y-q_A(\zeta)).
\end{equation}
The common physical root displacement has scale $N^{-1/4}\zeta$.
\end{theorem}
\begin{proof}
The finite set of clock probabilities is strictly positive on a two-sided neighbourhood of the baseline. Its derivatives through order three are bounded on a smaller neighbourhood. At $u=a/\sqrt N$, $v=h/\sqrt N$, expand each cell log probability. Normalization makes the mean linear term vanish, the expectation of the Hessian is minus the score covariance, and the summed cubic remainder is $O(N^{-1/2})$ on compact index sets. The centred Hessian fluctuation is $o_P(1)$. The finite-dimensional independent-group central limit theorem gives $\Delta_N\Rightarrow N(0,I)$ at the baseline and proves \eqref{eq:LAN}. Under the local alternatives its mean tends to $I(a,h)$ and its covariance to $I$, either by the same cell calculation or the likelihood change of measure. The efficient score in \eqref{eq:efficient-Y} has mean $Qh$ and covariance $Q$, proving the assertion after multiplication by $R$.

In the limiting score experiment the nuisance block is independent of the efficient block. Its mean is $I_{uu}a+I_{uv}h$, which is an unrestricted translation as $a$ varies. Discarding this nuisance translation yields \eqref{eq:efficient-Y}; this is a quotient, not a claim that an unknown nuisance parameter can be simulated without information. Profiling the efficient Gaussian likelihood over $h=q_A(\zeta)$ and completing the square gives \eqref{eq:profile-limit}. The map $q_A$ is proper for full-rank $A$, so the infimum is attained and its optimizing $\zeta$ are bounded whenever $Y$ is bounded. Uniform LAN therefore also gives the iterated compact-localization limit of profiled likelihoods: first restrict $(a,\zeta)$ to a fixed compact set, pass to the limit, and then exhaust the limiting experiment by such sets. No assertion about remote maximizers of an unrestricted finite-sample likelihood is required. Finally $q_A(N^{-1/4}\zeta)=N^{-1/2}q_A(\zeta)$, proving the displacement scale.
\end{proof}
The unrestricted vector $Y$ is a statistic in the limiting experiment, not the constrained maximum likelihood estimator. The latter is a metric projection onto $\R^k_{\ge0}$, or onto $q_A(\R^d)$ in the synchronized experiment. At the boundary it need not be normal; for the nonconvex synchronized image the projection can be set-valued. The distance value in \eqref{eq:profile-limit} is nevertheless unambiguous.

\begin{corollary}[Sharp root-amplitude scale on a physical ray]\label{cor:root-rate}
Fix $a_0\in\R^d\setminus\{0\}$ and $v_0=q_A(a_0)\ne0$. On the one-sided submodel with physical displacement $t a_0$, $0\le t\le T N^{-1/4}$, and nuisance $u=a/\sqrt N$ in a fixed bounded set, the minimax squared risk for estimating $t$ is of order $N^{-1/2}$, for each fixed $T>0$. Thus the amplitude scale is $N^{-1/4}$, whereas its squared amplitude is on the $N^{-1/2}$ parameter scale.
\end{corollary}
\begin{proof}
The variance parameter is $v=t^2v_0$. Set $J_0=v_0^{\mathsf T}Qv_0>0$ and use the baseline efficient score to define
\[
 \widehat b_N=\frac{v_0^{\mathsf T}QY_N}{\sqrt N J_0},\qquad
 \widehat t_N=\sqrt{\max\{\widehat b_N,0\}}.
\]
For completeness, let $s_{je}$ be the baseline score of cell $e$ in group $j$. There are finitely many such bounded vectors. Taylor expansion of the cell probabilities, rather than convergence in distribution, gives
\[
 \sup_{(a,h)\in K}\|\mathbb E_{a,h}\Delta_N-I_N(a,h)\|
       \le C_KN^{-1/2},\qquad
 \sup_{(a,h)\in K}\|\operatorname{Var}_{a,h}\Delta_N\|_{\rm op}
       \le C_K,
\]
where $I_N$ uses the actual allocations $n_j/N$ and $I_N\to I$. The first bound sums the bounded second-order probability remainders; the second sums variances of independent bounded scores. The efficient linear map consequently has mean $h+o(1)$ and bounded covariance, uniformly on $K$. With $h=\sqrt N t^2v_0$, the definition of $\widehat b_N$ now gives bias $o(N^{-1/2})$ and variance $O(N^{-1})$, and hence
$\mathbb E(\widehat b_N-t^2)^2\le C/N$.
This argument supplies the moment control separately from LAN. For $b\ge0$,
$(\sqrt{\max\{x,0\}}-\sqrt b)^2\le|x-b|$.
Cauchy--Schwarz therefore gives the upper bound $C'/\sqrt N$.

For the lower bound take the two physical alternatives $t=0$ and $t=\tau N^{-1/4}$ with the same nuisance, where $0<\tau\le T$ is sufficiently small. Their likelihood divergence is bounded by a constant times $\tau^4$, by \eqref{eq:LAN} or its finite-cell Taylor proof. Their total variation is consequently bounded away from one. The two-point testing bound for squared loss gives a positive constant times the squared separation $\tau^2N^{-1/2}$. The claim concerns this labelled nonnegative ray; it does not assert identifiability of the sign of a general loading vector.
\end{proof}

\section{Direct count-to-contact experiment comparison}\label{sec:score-contact}
Here the observations differ from fixed-count categorical marks. Let the mark laws, hence $\kappa_j$, be known, and observe independent exposures
\[
 N_j\sim\operatorname{Pois}(N\lambda_j),\qquad j=1,\ldots,m.
\]
This is the known-mark-law subexperiment of an observed-exposure protocol. Fix an interior centre $\lambda^0>0$. Write $D$ for the derivative of the native moment functional at that centre:
\[
 D_{\cdot j}=-\frac{\rho_j}{(\lambda_j^0)^2\kappa_j}w_{2n}(T_j),
 \quad \Lambda=\diag(\lambda_j^0),\quad
 \Sigma=D\Lambda D^{\mathsf T}>0.
\]
The positivity follows from Vandermonde rank. For a local moment index $h\in\R^p$ in a fixed compact ball, define the least-favourable target submodel
\begin{equation}\label{eq:least-favourable}
 \lambda_N(h)=\lambda^0+N^{-1/2}B_*h,
 \qquad B_*=\Lambda D^{\mathsf T}\Sigma^{-1}.
\end{equation}
It has $\theta(\lambda_N(h))=\theta(\lambda^0)+h/\sqrt N+O(N^{-1})$. Its local information is $\Sigma^{-1}$, since $DB_*=I_p$ and $B_*^{\mathsf T}\Lambda^{-1}B_*=\Sigma^{-1}$.

Let $\LL$ be the known normal-compression operator, possibly rank deficient, and put $r=\rank\LL$. Choose an orthonormal coordinate frame $F$ for its range. Generate independent $U_j\sim\operatorname{Unif}(-1/2,1/2)$, independently of the counts, and retain only
\begin{equation}\label{eq:direct-statistic}
 T_N=F^{\mathsf T}\LL D\frac{N_\cdot-N\lambda^0+U}{\sqrt N}.
\end{equation}
The randomization law is part of the protocol; its realized auxiliary values are not retained. This statistic is formed directly by fixed linear weights on the counts. No estimator of $\theta$ is constructed first. The small randomization prevents lattice encoding in this protocol: without such randomization or a genuine finite-precision coarsening, an exact real-valued linear encoding of an integer vector can be injective even when its real matrix has deficient rank. Weak convergence alone would not exclude that lattice encoding.

For experiments $\mathcal A=(P_h)$ and $\mathcal B=(Q_h)$ on the same parameter set, use
\[
 \delta(\mathcal A,\mathcal B)=\inf_K\sup_h\|KP_h-Q_h\|_{\rm TV},
 \qquad \Delta(\mathcal A,\mathcal B)
       =\max\{\delta(\mathcal A,\mathcal B),\delta(\mathcal B,\mathcal A)\}.
\]
The infimum is over parameter-independent Markov kernels and
$\|P-Q\|_{\rm TV}=\sup_A|P(A)-Q(A)|$, equal to one half the $L^1$ distance for densities. All comparisons below use a fixed compact local parameter set.

\begin{lemma}[Quantitative jittered Poisson comparison]\label{lem:poisson-TV}
Uniformly for $h$ in a fixed compact set, the distribution of
\[
 X_N^\circ=(N_\cdot-N\lambda^0+U)/\sqrt N
\]
under \eqref{eq:least-favourable} is within $C N^{-1/5}$ in total variation of
$N(B_*h,\Lambda)$. The constant depends on the fixed dimension, positive exposure bounds, and the compact index set. The full jittered vector and the raw counts are exactly equivalent experiments, by independent jittering and rounding. The raw local target experiment is therefore asymptotically equivalent to $Z\sim N(h,\Sigma)$.
\end{lemma}
\begin{proof}
We give the density estimate, including the lattice interpolation and tails. In one coordinate write $K\sim\operatorname{Pois}(N\lambda)$, with $0<\lambda_-\le\lambda\le\lambda_+$, and centre first at its own mean. The density $f_{N,\lambda}$ of $(K-N\lambda+U)/\sqrt N$ is constant, equal to $\sqrt N\Pr(K=j)$, on each cell of width $N^{-1/2}$. Put $z_j=(j-N\lambda)/\sqrt N$ and let $\varphi_\lambda$ be the centred normal density with variance $\lambda$. Stirling's formula with remainder
\[
 \log j!=(j+\tfrac12)\log j-j+\tfrac12\log(2\pi)+r_j,
       \qquad 0<r_j<1/(12j),
\]
and the Taylor formula for $\log(1+u)$ give, uniformly for
$|z_j|\le 2N^{1/10}$ and all sufficiently large $N$,
\[
 \left|\log\frac{\sqrt N\Pr(K=j)}{\varphi_\lambda(z_j)}\right|
 \le C\{N^{-1/2}(1+|z_j|^3)+N^{-1}(1+|z_j|^4)\}.
\]
Indeed $j=N\lambda(1+z_j/(\sqrt N\lambda))$, the quadratic term is $-z_j^2/(2\lambda)$, and the remaining logarithmic, cubic, and Stirling terms have the displayed bounds. Within a cell,
$|\log\varphi_\lambda(z)-\log\varphi_\lambda(z_j)|
 \le C N^{-1/2}(1+|z_j|)$.
Consequently on $|z|\le R_N=N^{1/10}$ the density ratio differs from one by at most $C N^{-1/5}$. Integrating this relative bound costs at most $C N^{-1/5}$. Outside this region,
\[
 \mathbb E\{(K-N\lambda+U)/\sqrt N\}^2
           =\lambda+(12N)^{-1},
\]
so Chebyshev's inequality bounds the step-density tail by $C R_N^{-2}$; the same bound holds for the Gaussian tail. Thus
\begin{equation}\label{eq:poisson-quantitative}
 \|f_{N,\lambda}-\varphi_\lambda\|_{L^1}
                       \le C N^{-1/5}.
\end{equation}
Translation to the centre $N\lambda^0$ changes neither the bound nor its uniformity. The finitely many small $N$ are absorbed by increasing $C$.

For $m$ independent coordinates, the total variation distance between products is at most the sum of the coordinate distances. This proves the same order for covariance $\operatorname{diag}\lambda_N(h)$. To replace it by $\Lambda$, differentiate the normal density with respect to each variance on the fixed positive compact interval. The $L^1$ norms of these derivatives are uniformly bounded (they are the normal expectations of the absolute variance scores). Since $\|\lambda_N(h)-\lambda^0\|=O(N^{-1/2})$ uniformly, covariance replacement costs $O(N^{-1/2})$. This proves the asserted quantitative comparison.

Rounding $\sqrt N X_N^\circ+N\lambda^0$ to the nearest integer recovers each count almost surely, while independent jittering gives the reverse Markov kernel. On arbitrary Gaussian inputs the rounding kernel clips negative integers to zero, so it is defined on the entire observation space. In the limiting normal vector $X\sim N(B_*h,\Lambda)$, the statistic $Z=DX$ has law $N(h,\Sigma)$ and is sufficient: conditional simulation is
\[
 X=B_*Z+E,\qquad
 E\sim N(0,\Lambda-\Lambda D^{\mathsf T}\Sigma^{-1}D\Lambda),
\]
with $E$ independent and its law independent of $h$. These kernels, the total variation approximation, and contraction of total variation under kernels prove the experiment comparison uniformly on the compact parameter set. In particular this proves the asserted local Le Cam comparison, not merely a central limit theorem.
\end{proof}

\begin{theorem}[Known-mark local information retained by contact scores]\label{thm:contact-experiment}
With known mark laws and the preassigned centre and target submodel \eqref{eq:least-favourable}, the experiment retaining \eqref{eq:direct-statistic} is, uniformly on compact local parameter balls, asymptotically equivalent to
\begin{equation}\label{eq:contact-limit}
 T\sim N(F^{\mathsf T}\LL h,\ F^{\mathsf T}\LL\Sigma\LL^{\mathsf T}F).
\end{equation}
Its limiting information matrix for $h$ is
\begin{equation}\label{eq:contact-information}
 I_{\rm contact}=\LL^{\mathsf T}(\LL\Sigma\LL^{\mathsf T})^\dagger\LL
 =\Sigma^{-1/2}P_{\ran(\Sigma^{1/2}\LL^{\mathsf T})}\Sigma^{-1/2}.
\end{equation}
It has rank $r$ and kernel $\ker\LL$. The contact-score reduction is locally asymptotically equivalent to the raw target subexperiment if and only if $\LL$ is injective. Hence it retains all local moment information exactly when the within-contact polynomial products span $V_{2n}$.
\end{theorem}
\begin{proof}
Apply the fixed linear map $F^{\mathsf T}\LL D$ to the total variation approximation in Lemma~\ref{lem:poisson-TV}. Contraction gives convergence in total variation to \eqref{eq:contact-limit}, whose covariance is positive definite in its $r$ coordinates. Identity kernels therefore compare these experiments in both directions with error tending to zero. The usual normal-shift likelihood gives the first expression in \eqref{eq:contact-information}. For $M=\LL\Sigma^{1/2}$, the identity $M^{\mathsf T}(MM^{\mathsf T})^\dagger M=P_{\ran M^{\mathsf T}}$ gives the second. Its rank and kernel follow immediately.

If $\LL$ is injective, $Z$ is recovered from $\LL Z$ by a left inverse. Thus \eqref{eq:contact-limit} is equivalent to $N(h,\Sigma)$, and Lemma~\ref{lem:poisson-TV} gives equivalence to the raw target subexperiment. Conversely choose distinct $h_0,h_1$ in a local parameter ball with $h_1-h_0\in\ker\LL$. Their limiting contact distributions coincide, whereas $N(h_0,\Sigma)$ and $N(h_1,\Sigma)$ have strictly positive total variation distance. By the triangle inequality, any kernel purporting to reconstruct both raw limit laws has error at least half that distance. The same positive lower obstruction holds asymptotically by the uniform comparisons. The polynomial criterion is Theorem~\ref{thm:fibre}.
\end{proof}
The uncentred contact block is $B_\nu(\theta)=\LL_\nu\theta$. Its local fluctuation is precisely the linear type appearing in \eqref{eq:direct-statistic}. Passing to inverse-Hessian contact scores applies the fixed invertible block transformation $E_\nu\mapsto-B_{\nu,0}^{-1}E_\nu B_{\nu,0}^{-1}$, so it preserves the information statement. These are correlated raw-data score statistics, not independently noisy observations of distance values.

The centre $\lambda^0$ and the local submodel \eqref{eq:least-favourable} are specified before observing the data. The theorem does not claim global sufficiency when both exposure and mark laws are unknown, or that a plug-in estimate of the centre automatically gives the same comparison. It supplies a genuine direct reduction and a complete local information-loss calculation in the stated full-dimensional target experiment. When $\LL$ is injective, exact range-compatible score encoding does not introduce a $\gamma^{-2}$ statistical variance penalty; independent measurement errors and numerical perturbations have the distinct stability behaviour of Theorem~\ref{thm:stability} and Appendix~\ref{app:offset}.


\subsection{Deterministic finite precision and the full exposure tangent model}
The auxiliary randomization above prevents an exact real-valued encoding of the integer count vector. A finite-precision observation provides another, deterministic, way to remove that encoding.
\begin{theorem}[Deterministic quantized score reduction]\label{thm:quantized-score}
Keep the known-mark target submodel and put $C=F^{\mathsf T}\LL D$. Suppose $r\ge1$. For a deterministic mesh $\eta_N>0$, retain only
\[
 Q_N=\left\lfloor \eta_N^{-1}
          C(N_\cdot-N\lambda^0)/\sqrt N\right\rfloor\in\mathbb Z^r,
\]
with the floor taken coordinatewise. If $\eta_N\to0$ and
$\sqrt N\eta_N\to\infty$, this experiment is asymptotically equivalent to \eqref{eq:contact-limit}. More precisely, for $0<\eta_N\le1$,
\begin{equation}\label{eq:quantized-bound}
 \Delta(\mathcal Q_N,\mathcal T)
 \le C_0\{N^{-1/5}+(\sqrt N\eta_N)^{-1}+\eta_N\}.
\end{equation}
In particular $\eta_N=N^{-1/4}$ is admissible. The information and equivalence criterion of Theorem~\ref{thm:contact-experiment} therefore hold for this deterministic protocol as well.
\end{theorem}
\begin{proof}
Couple the unjittered score $S_N=C(N_\cdot-N\lambda^0)/\sqrt N$ to $S_N^\circ=CX_N^\circ$ with the same counts. Their coordinatewise difference is at most $d_N=C_1/\sqrt N$, deterministically. The two quantizers can differ only if a coordinate of $S_N^\circ$ lies within $d_N$ of the grid $\eta_N\mathbb Z$. By Lemma~\ref{lem:poisson-TV}, its joint law is within $C N^{-1/5}$ of the nonsingular normal law in \eqref{eq:contact-limit}.

For a univariate normal density $f$ with variance in a fixed positive interval, the grid-boundary probability is at most $C d_N/\eta_N$ if $d_N\le\eta_N/2$. To see this uniformly in its mean, sum the supremum of $f$ on each grid interval and use
\[
 \eta_N\sum_{j\in\mathbb Z}\sup_{[j\eta_N,(j+1)\eta_N]}f
       \le 1+\eta_N\int_{\R}|f'(x)|\,dx.
\]
There are at most two boundary pieces per interval. A union bound over the $r$ coordinates proves
$\Pr\{\lfloor S_N/\eta_N\rfloor\ne
       \lfloor S_N^\circ/\eta_N\rfloor\}
 \le C\{N^{-1/5}+d_N/\eta_N\}$.
When $d_N>\eta_N/2$ the asserted bound follows by increasing its constant. Quantization contracts total variation, so $Q_N$ is within the same bound of a quantized Gaussian, uniformly in the local index.

The kernel adding an independent uniform point in the reported grid cell reconstructs the cell-average Gaussian density. Its $L^1$ distance from the Gaussian is at most
$C_r\eta_N\int\|\nabla f(x)\|\,dx\le C\eta_N$,
by the integral mean-value inequality on cubes. The reverse kernel is deterministic quantization. Combining these kernels proves both deficiencies in \eqref{eq:quantized-bound}. Randomization is used only by this comparison kernel, not by the observed statistic $Q_N$. The case $r=0$ is the one-point experiment and needs no mesh.
\end{proof}
An actual finite alphabet may be obtained by an overflow symbol outside a ball of radius $N^{1/10}$. Uniform second-moment bounds add only $O(N^{-1/5})$ to the comparison. Thus infinite digital precision is not hidden in the conclusion.

The least-favourable submodel has a precise relationship to the full exposure tangent experiment, rather than being an intrinsic physical distance observation.
\begin{proposition}[Full exposure model and nuisance quotient]\label{prop:full-exposure}
For known marks and arbitrary bounded exposure perturbations $g\in\R^m$, set
$\lambda_N(g)=\lambda^0+g/\sqrt N$. The raw count experiment converges in Le Cam distance, uniformly on compact sets, to $X\sim N(g,\Lambda)$. In this Gaussian limit the transformation
\[
 Z=DX,\qquad E=(I_m-B_*D)X
\]
has independent components, with means
$h=Dg$ and $\xi=(I_m-B_*D)g\in\ker D$. Quotienting the limit experiment by translations of this unrestricted nuisance gives exactly $Z\sim N(h,\Sigma)$; its contact compression has information \eqref{eq:contact-information}. The least-favourable submodel is the section $\xi=0$.
\end{proposition}
\begin{proof}
The density proof of Lemma~\ref{lem:poisson-TV} applies unchanged to $g$ in a compact set; jittering and rounding give both kernels. Since $DB_*=I_p$, $g=B_*h+\xi$ is a bijective linear parametrization by $h\in\R^p$ and $\xi\in\ker D$. The cross-covariance is
\[
 D\Lambda(I_m-B_*D)^{\mathsf T}
       =D\Lambda-\Sigma B_*^{\mathsf T}=0.
\]
The Gaussian components are therefore independent, the covariance of $Z$ is $\Sigma$, and the law of $E-\xi$ does not depend on either parameter. Translation by $\ker D$ acts transitively on its nuisance means and leaves $Z$ fixed. This proves the quotient assertion. It does not supply a kernel simulating an unknown $\xi$ from $Z$, nor assert equivalence to the full experiment with that nuisance retained. The compact-set comparisons are made before this unrestricted Gaussian translation quotient; a compact nuisance set is not silently replaced by a translation group.
\end{proof}

The centre remains an experimental input. A simple sufficient condition for replacing it by an independent pilot estimate is
$\sqrt N\|\widehat\lambda^0-\lambda^0\|\to0$ in probability, together with a fixed positive compact range for the pilot and continuity of the weights. On events where the scaled error is small, the centring shift is $o(1)$, the weights and Gaussian covariances converge, and the same uniform density bounds and mixture kernels apply. A root-$M$ pilot gives this sufficient condition when $M/N\to\infty$. This is a sufficient, not an optimal, pilot regime. An ordinary pilot of size comparable to $N$ does not meet it: its random centring error must then be included in the limit experiment. Nothing here asserts automatic plug-in sufficiency for unknown marks or an unrecorded pilot.
